\section{Discussion}
The evidence does not support a single hierarchy of VR methods. It supports a hierarchy of claims. A narrow, controlled comparison can justify a regime effect; a heterogeneous collection may justify a taxonomy or moderator hypothesis; neither automatically supports a universal ranking. This is why semantic compatibility must precede aggregation.

The RDW results make the point experimentally. S2C is attractive when simplicity and low decision cost dominate, but the difficult G03 geometry changes both completion and reset-burden ordering. Vis-Poly pays a much larger computational price yet becomes comparatively robust in that geometry. P2R occupies a different part of the trade space. The scientifically interesting object is therefore the boundary at which method ordering changes, not a single global best controller.

Prior art constrains the novelty claim. Long-standing RDW taxonomies, simulation studies, and recent reviews already document geometry, path, prediction, reset, and controller-family variation \cite{Nilsson2018FifteenYears,Hodgson2013FourApproaches,Azmandian2015PhysicalSpace,Prinz2023Taxonomies}. Context-aware, learned, predictive, and geometry-aware redirection already exist \cite{Azmandian2022Adaptive,Strauss2020RL,Chen2024APFS2T,Wang2026DGMRDW}. We therefore do not present a new adaptive controller. A defensible future problem is narrower: jointly predict failure probability, burden conditional on completion, decision cost, and uncertainty, and validate the policy on held-out geometries.

The same logic generalizes beyond locomotion. Presence and embodiment should not be collapsed across instruments; cybersickness should be interpreted with exposure and susceptibility; education and health require comparator equivalence; rendering and networking should be evaluated end to end; security needs practical validation; and accessibility needs to be treated as a design obligation rather than a post hoc attribute \cite{Caserman2021Cybersickness,Merchant2014VREducationMeta,Woon2021NursingVR,Wang2023FoveatedSurvey,Giaretta2025VRSecurity,Dudley2023InclusiveImmersion}.

\section{Threats to Validity}
The evidence-library pipeline is curated rather than reconstructible from raw per-database discovery exports. It reliably accounts for the registered 174-record library and subsequent screening, but early database-specific yield counts were not preserved and are not retroactively inferred. Two reports were not retrieved and therefore were not treated as full-text assessed evidence.

The broad review is heterogeneous by design; the comparability framework reduces invalid synthesis but cannot eliminate publication bias, incomplete reporting, or hidden dependence among study families. The RDW benchmark is simulation-based, uses four selected geometries, and evaluates three executable controllers under a common implementation environment. Simulation is useful for controlled RDW analysis but still requires explicit validation against live-user behavior \cite{Hodgson2013FourApproaches,Azmandian2022SimulationValidation,Hirt2022Chaotic}. Its internal regime effect is controlled, but live-user and broader-geometry validation remain necessary. The computational-cost benchmark is also hardware/software dependent and measures decision logic rather than rendering, tracking, reset execution, or network latency.

\section{Reproducibility and Data/Code Availability}
The study was developed in the sequence represented in Algorithm~\ref{alg:pipeline}: first the evidence/comparability concept, then the executable benchmark software, then versioned workflow execution, then artifact-level audit and statistics. The public \repo{} repository serves as the archival code, data, and provenance record for this research. It contains the benchmark implementation, test suite, machine-readable review registers, corrected N07--N09 quantitative outputs, and provenance notes needed to reproduce the reported tables and figures. The repository is not treated as a manuscript-generation mechanism.

The final quantitative claims use the corrected N07--N09 artifacts. Earlier pre-forensic summaries are retained only as provenance and are superseded where they disagree with immutable workflow artifacts. The final package also includes the exact source CSVs used for the tables and figures and a deterministic figure-generation script.

\section{Conclusion}
VR results become more useful when the conditions under which they hold are treated as part of the result itself. Across the reviewed domains, apparently conflicting conclusions often reduce to differences in task, construct, population, hardware, intervention, metric, or validation level. In the pre-specified contradiction audit, seven of eight families were comparability-limited, while the controlled RDW family exposed a genuine regime effect.

The benchmark makes that conclusion concrete. Visibility-Polygon, S2C, and P2R differed in completion, reset burden, and decision cost, and their ordering changed with geometry. No controller was universally best. We therefore reject an artificial controller-novelty claim and retain the stronger, more defensible result: VR methods should be compared under explicit operating regimes, with failures separated from conditional performance and with raw execution evidence preserved through analysis.

The broader methodological implication is equally important. Comparability should precede ranking; reproducibility claims should identify the level actually demonstrated; negative outcomes should remain visible; and research gaps should progress from observation to corroboration, testability, benchmark readiness, and experimental evaluation. Under this view, the most productive open problems are not unsupported absences in the literature but reproducible boundaries where existing methods change behavior, fail, or cease to generalize.

\section*{Data and Code Availability}
The source code, benchmark configuration, corrected analysis tables, statistical outputs, and evidence registers are archived in \repo{}. The submission package contains the exact figure source data and code excerpts required to reproduce the manuscript figures and quantitative summaries.
